% Chapter C: Probability and Statistics Review
% A First Course in Causal Inference - Peng Ding

\chapter{概率论与统计学复习}

\section{本章导论}

本书假设读者具备概率论与统计学的基础知识。然而，这些基础知识在因果推断中扮演着至关重要的角色——我们即将遇到的许多概率计算，都依赖于本章所复习的核心概念。

\textbf{为什么需要概率论复习？}
在后续章节中，我们将频繁处理以下问题：
\begin{itemize}
  \item 计算潜在结果的期望与方差
  \item 理解条件分布与条件期望
  \item 运用中心极限定理建立渐近分布
  \item 使用 Delta 方法处理非线性函数的渐近性质
\end{itemize}

这些问题看似基础，但它们构成了整个因果推断理论大厦的基石。正如我们在第一章中看到的，Yule-Simpson 悖论的本质就是隐藏在概率计算背后的非线性效应。

\textbf{本章结构：}
本章按以下线索展开。首先复习概率论的核心工具：期望、方差、协方差及其运算性质（tower property、Eve's law）。接着介绍多元正态分布与常用小样本分布。然后讨论统计学核心概念：点估计、区间估计、假设检验。最后以两个经典统计问题（Behrens-Fisher 问题与 Fieller-Creasy 问题）作为延伸阅读。

\section{概率基础回顾}

\subsection{期望、方差与协方差}

在因果推断中，我们最常处理的量是随机变量的期望与方差。给定随机变量 $Y$，其期望定义为
\[
\mathbb{E}(Y) = \int y \, dF(y),
\]
方差定义为
\[
\var(Y) = \mathbb{E}\{(Y - \mathbb{E}Y)^2\} = \mathbb{E}(Y^2) - \{\mathbb{E}(Y)\}^2 .
\]

\textbf{关键性质——Tower Property（Adam's Law）：} 对于任意随机变量 $A$ 和 $B$，有
\begin{equation}\label{eq:tower-property}
\mathbb{E}(A) = \mathbb{E}\{\mathbb{E}(A \mid B)\}.
\end{equation}

这个看似简单的性质，却是条件期望理论的起点。它的名字来源于嵌套期望的结构：外层期望作用于内层条件期望。

\textbf{方差分解——Eve's Law：} 若读者曾在研究生阶段学习统计，可能会记得 Joe Blitzstein 教授将下列公式称为"Eve's Law"，因为公式中出现了"EVVE"的结构：
\begin{equation}\label{eq:eves-law}
\var(A) = \mathbb{E}\{\var(A \mid B)\} + \var\{\mathbb{E}(A \mid B)\}.
\end{equation}

这个公式告诉我们：一个量的总方差可以分解为条件方差均值加上条件期望的方差。在因果推断中，这个分解是理解混杂与选择偏倚的核心工具。

\textbf{协方差与相关系数：} 两个随机变量 $Y$ 和 $X$ 的协方差定义为
\[
\cov(Y, X) = \mathbb{E}\{(Y - \mathbb{E}Y)(X - \mathbb{E}X)\} = \mathbb{E}(YX) - \mathbb{E}(Y)\mathbb{E}(X).
\]

Pearson 相关系数定义为
\begin{equation}\label{eq:pearson-correlation}
\rho_{YX} = \frac{\cov(Y, X)}{\sqrt{\var(Y)\var(X)}} .
\end{equation}

相关系数衡量的是 $Y$ 与 $X$ 之间的线性依赖程度，其取值在 $[-1, 1]$ 之间。

\textbf{多元相关系数：} 若 $Y$ 是随机变量，$X$ 是随机向量，则多元相关系数定义为
\[
R_{YX}^2 = \corr^2(Y, X) = \frac{\cov(Y, X) \, \cov(X)^{-1} \, \cov(X, Y)}{\var(Y)} .
\]
它衡量 $Y$ 与 $X$ 向量之间线性依赖的程度，但这个定义在 $Y$ 和 $X$ 之间不对称。

\textbf{协方差的分解：} 类似于方差分解，我们也有协方差的分解公式（见 \cref{sec:derivation-covariance-decomposition}）：
\[
\cov(A_1, A_2) = \mathbb{E}\{\cov(A_1, A_2 \mid B)\} + \cov\{\mathbb{E}(A_1 \mid B), \mathbb{E}(A_2 \mid B)\}.
\tag*{\footnote{推导见附录 \cref{sec:derivation-covariance-decomposition}}}
\]

这个公式在因果推断中非常重要——它帮助我们理解在控制某个变量 $B$ 后，$A_1$ 与 $A_2$ 之间的关系如何变化。

\subsection{Cauchy-Schwarz 不等式}

Cauchy-Schwarz 不等式是概率论中最基本的不等式之一。对于任意两个随机变量 $A$ 和 $B$，有
\begin{equation}\label{eq:cauchy-schwarz}
|\mathbb{E}(AB)| \leq \sqrt{\mathbb{E}(A^2)\mathbb{E}(B^2)} .
\end{equation}

等号成立的充要条件是存在常数 $\beta$ 使得 $B = \beta A$。

将随机变量中心化（即减去均值）后，可得协方差形式：
\[
|\cov(A, B)| \leq \sqrt{\var(A)\var(B)} .
\]

这正是 \cref{eq:pearson-correlation} 中相关系数绝对值不超过 $1$ 的根本原因。

\section{多元正态分布}

在因果推断中，多元正态分布是一类核心的分布假设。许多理论结果都是在多元正态的框架下建立的。

设随机向量 $(Y_1, Y_2)^\top$ 服从二元正态分布：
\[
\begin{pmatrix}
Y_1 \\ Y_2
\end{pmatrix}
\sim N\left(
\begin{pmatrix}
\mu_1 \\ \mu_2
\end{pmatrix},
\begin{pmatrix}
\sigma_1^2 & \sigma_{12} \\
\sigma_{12} & \sigma_2^2
\end{pmatrix}
\right) .
\]

\textbf{边缘分布：} 边际分布仍是正态：
\[
Y_1 \sim N(\mu_1, \sigma_1^2), \quad Y_2 \sim N(\mu_2, \sigma_2^2).
\]

\textbf{条件分布：} 若 $\sigma_2^2 > 0$，条件分布也是正态：
\begin{equation}\label{eq:conditional-normal}
Y_1 \mid Y_2 = y_2 \sim N\left(\mu_1 + \frac{\sigma_{12}}{\sigma_2^2}(y_2 - \mu_2), \sigma_1^2 - \frac{\sigma_{12}^2}{\sigma_2^2}\right).
\end{equation}

\cref{eq:conditional-normal} 揭示了多元正态的一个优雅性质：条件期望是线性函数，条件方差与条件均值无关。这个性质在偏相关分析中至关重要。

\textbf{偏相关系数：} 在第一章中我们遇到过偏相关系数。对于三元正态 $(Y, X, Z)^\top$，$Y$ 和 $Z$ 的偏相关系数定义为条件分布 $(Y, Z) \mid X$ 中的相关系数：
\begin{equation}\label{eq:partial-correlation}
\rho_{YZ|X} = \frac{\rho_{YZ} - \rho_{YX}\rho_{ZX}}{\sqrt{1-\rho_{YX}^2}\sqrt{1-\rho_{ZX}^2}} .
\end{equation}

即使 $\rho_{YZ} > 0$，偏相关系数 $\rho_{YZ|X}$ 仍可能为负——这正是 Yule-Simpson 悖论在正态情形下的数学本质。

\section{常用小样本分布}

在统计学的小样本推断中，我们经常需要处理以下分布。

\subsection{$\chi^2$ 分布}

若 $X_1, \ldots, X_n \stackrel{\text{IID}}{\sim} N(0, 1)$，则
\[
\sum_{i=1}^n X_i^2 \sim \chi_n^2 ,
\]
其中 $\chi_n^2$ 表示自由度为 $n$ 的卡方分布。$\chi_n^2$ 分布的均值是 $n$，方差是 $2n$。

卡方分布在方差估计与拟合优度检验中扮演核心角色。

\subsection{$t$ 分布}

设 $X \sim N(0, 1)$，$Q_n \sim \chi_n^2$，且 $X \perp Q_n$，则
\[
\frac{X}{\sqrt{Q_n / n}} \sim t_n ,
\]
其中 $t_n$ 表示自由度为 $n$ 的 $t$ 分布。当 $n=1$ 时，$t_1$ 分布就是著名的 Cauchy 分布——它没有有限的均值。

$t$ 分布在小样本情形下用于构造均值估计的置信区间。

\section{极限定理}

\subsection{依概率收敛与依分布收敛}

在建立渐近理论之前，我们需要明确两种基本的收敛概念。

\begin{Definition}[依概率收敛]\label{def:convergence-probability}
若对任意 $\varepsilon > 0$，有
\[
\Pr(|X_n - X| > \varepsilon) \rightarrow 0 \quad \text{as } n \to \infty,
\]
则称序列 $(X_n)_{n \geq 1}$ 依概率收敛到 $X$，记作 $X_n \xrightarrow{p} X$。
\end{Definition}

\begin{Definition}[依分布收敛]\label{def:convergence-distribution}
若对 $\Pr(X \leq x)$ 的所有连续点 $x$，有
\[
\Pr(X_n \leq x) \rightarrow \Pr(X \leq x) \quad \text{as } n \to \infty,
\]
则称序列 $(X_n)_{n \geq 1}$ 依分布收敛到 $X$，记作 $X_n \xrightarrow{d} X$。
\end{Definition}

依概率收敛强于依分布收敛。

\subsection{大数定律}

\begin{Theorem}[弱大数定律]\label{eq:lln}
若 $X_1, \ldots, X_n \stackrel{\text{IID}}{\sim} X$ 且 $\mathbb{E}|X| < \infty$，则样本均值依概率收敛到总体均值：
\[
\bar{X} = \frac{1}{n}\sum_{i=1}^n X_i \xrightarrow{p} \mathbb{E}(X).
\]
\end{Theorem}

大数定律告诉我们：随着样本量增大，样本均值越来越接近总体均值。这是统计推断的理论基础。

\textbf{历史注记：} 大数定律的研究可以追溯到 Bernoulli（1713），但 Bernoulli 证明的是现在所称的"弱大数定律"。真正严谨的证明由 Chebyshev（1867）和 Borel（1909）给出。

\subsection{中心极限定理}

\begin{Theorem}[中心极限定理]\label{eq:clt}
若 $X_1, \ldots, X_n \stackrel{\text{IID}}{\sim} X$ 且 $\var(X) < \infty$，则
\[
\frac{\bar{X} - \mathbb{E}(X)}{\sqrt{\var(X) / n}} \xrightarrow{d} N(0, 1).
\]
\end{Theorem}

中心极限定理是统计学中最重要的定理之一：无论总体分布是什么，只要样本量足够大，标准化后的样本均值就近似服从标准正态分布。这使得我们能够在不知道总体分布的情况下进行统计推断。

\textbf{历史注记：} 中心极限定理的发展经历了漫长历程。De Moivre（1733）首先证明二项分布的近似正态性，Laplace（1810）推广了这一结果。Lindeberg（1922）和 Feller（1935）给出了最一般的条件。

\section{Delta 方法}

Delta 方法是建立非线性函数渐近正态性的有力工具。

\begin{Theorem}[Delta 方法]\label{eq:delta-method}
设 $\sqrt{n}(X_n - \mu) \xrightarrow{d} N(0, \Sigma)$，函数 $g(x)$ 在 $\mu$ 处有非零导数 $\nabla g(\mu)$。则
\[
\sqrt{n}\{g(X_n) - g(\mu)\} \xrightarrow{d} N\left(0, (\nabla g(\mu))^\top \Sigma \, \nabla g(\mu)\right).
\]
\end{Theorem}

Delta 方法的本质是泰勒展开：
\[
g(X_n) - g(\mu) \approx (\nabla g(\mu))^\top (X_n - \mu).
\]
因此，$g(X_n)$ 的渐近分布等价于一个线性变换后的正态分布。

\textbf{应用示例——比率的渐近正态性：}\footnote{推导见附录 \cref{sec:derivation-ratio-asymptotic}} 设
\[
\sqrt{n}\begin{pmatrix} Y_n - \mu_Y \\ X_n - \mu_X \end{pmatrix} \xrightarrow{d} N\left(\begin{pmatrix} 0 \\ 0 \end{pmatrix}, \begin{pmatrix} \sigma_Y^2 & \sigma_{YX} \\ \sigma_{YX} & \sigma_X^2 \end{pmatrix}\right),
\]
其中 $\mu_X \neq 0$。应用 Delta 方法，可得
\[
\sqrt{n}\left(\frac{Y_n}{X_n} - \frac{\mu_Y}{\mu_X}\right) \xrightarrow{d} N\left(0, \frac{\sigma_Y^2}{\mu_X^2} + \frac{\mu_Y^2\sigma_X^2}{\mu_X^4} - \frac{2\mu_Y\sigma_{YX}}{\mu_X^3}\right).
\tag*{\footnote{推导见附录 \cref{sec:derivation-ratio-asymptotic}}}
\]

这个结果在因果推断中非常重要——许多估计量（如风险比、优势比）都可以表示为某种比率形式。

\section{统计学基础}

\subsection{点估计}

设 $\theta$ 为感兴趣参数，$\hat{\theta}$ 为基于数据的估计量。

\begin{Definition}[无偏性]\label{def:unbiasedness}
若对所有可能的 $\theta$ 和干扰参数 $\eta$，都有
\[
\mathbb{E}(\hat{\theta}) = \theta,
\]
则称 $\hat{\theta}$ 是 $\theta$ 的无偏估计。
\end{Definition}

\begin{Definition}[一致性]\label{def:consistency}
若当样本量趋于无穷时，对所有可能的 $\theta$ 和 $\eta$，有
\[
\hat{\theta} \xrightarrow{p} \theta,
\]
则称 $\hat{\theta}$ 是 $\theta$ 的一致估计。
\end{Definition}

无偏性不蕴含一致性，一致性也不蕴含无偏性。在大样本推断中，一致性通常是基本要求。

\subsection{置信区间}

\begin{Definition}[置信区间]\label{def:confidence-interval}
若对所有可能的 $\theta$ 和 $\eta$，都有
\[
\Pr(\hat{\theta}_L \leq \theta \leq \hat{\theta}_U) \geq 1 - \alpha,
\]
则称数据依赖区间 $[\hat{\theta}_L, \hat{\theta}_U]$ 为 $\theta$ 的覆盖率至少为 $1-\alpha$ 的置信区间。
\end{Definition}

\textbf{Wald 型置信区间：} 在许多问题中，我们先建立估计量的渐近正态性：
\[
\hat{\theta} \approx N(\theta, v),
\]
其中 $v$ 是渐近方差。然后用 $v$ 的估计 $\hat{v}$ 替换，得到 Wald 型置信区间：
\[
\hat{\theta} \pm z_{1-\alpha/2}\sqrt{\hat{v}} .
\]

\section{二乘二列联表推断}

在因果推断中，我们经常处理二值结果与二值处理的状态，此时数据可以用 $2 \times 2$ 列联表汇总。

\subsection{Fisher 精确检验}

设
\[
n_{11} \sim \text{Binomial}(n_1, p_1), \quad n_{01} \sim \text{Binomial}(n_0, p_0), \quad n_{11} \perp n_{01}.
\]

Fisher 证明了在原假设 $H_0: p_1 = p_0$ 下，条件分布 $n_{11} \mid n_{.1}$ 服从超几何分布：
\[
\Pr(n_{11} = k) = \frac{\binom{n_{.1}}{k}\binom{n - n_{.1}}{n_1 - k}}{\binom{n}{n_1}} .
\]

基于这个条件分布，我们可以计算精确的 $p$ 值——这就是 Fisher 精确检验的核心思想。

\textbf{动机：} 为什么 Fisher 要用条件分布？原因是：在原假设 $p_1 = p_0$ 下，边缘和 $n_{.1}$（第一列的总和）包含了关于共同成功率 $p$ 的信息，但条件分布 $n_{11} \mid n_{.1}$ 却不依赖这个未知参数——这正是精确推断的关键。

\subsection{风险差、风险比与优势比}

基于二项模型，我们可以估计各参数：

\begin{align}
\widehat{\text{RD}} &= \hat{p}_1 - \hat{p}_0, \label{eq:risk-diff} \\
\log\widehat{\text{RR}} &= \log\frac{\hat{p}_1}{\hat{p}_0}, \label{eq:log-risk-ratio} \\
\log\widehat{\text{OR}} &= \log\frac{\hat{p}_1/(1-\hat{p}_1)}{\hat{p}_0/(1-\hat{p}_0)} = \log\frac{n_{11}n_{00}}{n_{10}n_{01}} . \label{eq:log-odds-ratio}
\end{align}

渐近地，这些估计量都服从正态分布，其方差可由 Delta 方法得到。

\section{经典统计问题}

\subsection{Behrens-Fisher 问题}

Behrens-Fisher 问题关心的是：当两个总体的方差未知且可能不等时，如何构造两均值差的置信区间？这个问题没有精确解，只能求助于渐近方法或近似方法。

这个问题在因果推断中经常出现——当我们比较处理组与对照组的潜在结果均值时，往往不能假设两者方差相等。

\subsection{Fieller-Creasy 问题}

Fieller-Creasy 问题涉及比率参数的置信区间构造。例如，在因果推断中我们经常需要构造 $\mu_Y/\mu_X$ 的置信区间。这个问题的挑战在于：比率的分布本身不是正态的，需要特别处理。

这两个经典问题的共同特点是：它们都没有"标准答案"，需要根据具体问题选择合适的方法。这种灵活性在因果推断中同样常见。

\section{本章小结}

本章复习了因果推断所需的概率论与统计学基础知识：

\begin{enumerate}
  \item \textbf{期望、方差、协方差}：Tower Property 与 Eve's Law 是理解条件期望的核心工具
  \item \textbf{多元正态分布}：条件正态的线性结构是偏相关分析的基础
  \item \textbf{极限定理}：大数定律与中心极限定理是渐近推断的理论支柱
  \item \textbf{Delta 方法}：处理非线性函数渐近分布的标准工具
  \item \textbf{统计学基础}：无偏性、一致性、置信区间、假设检验
\end{enumerate}

这些内容将在后续章节中反复出现。建议读者在继续阅读之前，确保对这些概念有扎实的理解。

\section{习题}

\begin{Exercise}{\ref{ex:c.1} Tower Property 与 Eve's Law 证明}\label{ex:c.1}
证明 Tower Property $\mathbb{E}(A) = \mathbb{E}\{\mathbb{E}(A \mid B)\}$ 和 Eve's Law $\var(A) = \mathbb{E}\{\var(A \mid B)\} + \var\{\mathbb{E}(A \mid B)\}$。
\end{Exercise}

\begin{Exercise}{\ref{ex:c.2} Cauchy-Schwarz 不等式证明}\label{ex:c.2}
利用方差非负性，证明 $|\cov(A, B)| \leq \sqrt{\var(A)\var(B)}$。
\end{Exercise}

\begin{Exercise}{\ref{ex:c.3} 偏相关系数}\label{ex:c.3}
考虑三元正态随机向量 $(Y, X, Z)^\top \sim N(\mathbf{0}, \Sigma)$。证明偏相关系数公式（见 \cref{eq:partial-correlation}）。
\end{Exercise}

\begin{Exercise}{\ref{ex:c.4} Delta 方法应用}\label{ex:c.4}
设 $\sqrt{n}(X_n - \mu_X) \xrightarrow{d} N(0, \sigma_X^2)$，$\sqrt{n}(Y_n - \mu_Y) \xrightarrow{d} N(0, \sigma_Y^2)$，且 $X_n$ 与 $Y_n$ 渐近独立。证明 $X_n Y_n$ 的渐近正态性：
\[
\sqrt{n}(X_n Y_n - \mu_X \mu_Y) \xrightarrow{d} N(0, \mu_Y^2\sigma_X^2 + \mu_X^2\sigma_Y^2) .
\]
\end{Exercise}

\begin{Exercise}{\ref{ex:c.5} Fisher 精确检验}\label{ex:c.5}
写出 Fisher 精确检验的 $p$ 值计算公式，并说明其在何种场景下优于渐近检验。
\end{Exercise}

\begin{Exercise}{\ref{ex:c.6} 推荐阅读}\label{ex:c.6}
\begingroup
\parskip=0pt
\textbf{推荐阅读}：\cite{durrett2019} 是概率论的优秀教材，其中包含对极限理论的详尽讨论。\cite{van2000} 是渐近统计的标准参考。

\textbf{历史注记}：Eve's Law 最早由 Laplace 在 1774 年提出，但 Joe Blitzstein 教授在哈佛统计系任教期间广泛推广了这个名称。中心极限定理的发展历程同样精彩：De Moivre（1733）、Laplace（1810）、Lindeberg（1922）、Feller（1935）逐步建立了最一般的理论。
\endgroup
\end{Exercise}

% === 附录：推导 ===
\section{附录：推导证明}\label{sec:appendixC-derivations}

\subsection{协方差分解公式的推导}\label{sec:derivation-covariance-decomposition}

给定随机变量 $A_1, A_2, B$，我们有
\[
\cov(A_1, A_2) = \mathbb{E}(A_1 A_2) - \mathbb{E}(A_1)\mathbb{E}(A_2) .
\]

对条件期望应用 Tower Property：
\[
\mathbb{E}(A_1 A_2) = \mathbb{E}\{\mathbb{E}(A_1 A_2 \mid B)\},
\]
\[
\mathbb{E}(A_1) = \mathbb{E}\{\mathbb{E}(A_1 \mid B)\}, \quad \mathbb{E}(A_2) = \mathbb{E}\{\mathbb{E}(A_2 \mid B)\} .
\]

代入并整理，可得
\[
\cov(A_1, A_2) = \mathbb{E}\{\cov(A_1, A_2 \mid B)\} + \cov\{\mathbb{E}(A_1 \mid B), \mathbb{E}(A_2 \mid B)\} .
\tag*{\qed}
\]

\subsection{比率渐近正态性的推导}\label{sec:derivation-ratio-asymptotic}

设 $\sqrt{n}(Y_n - \mu_Y) \xrightarrow{d} N(0, \sigma_Y^2)$，$\sqrt{n}(X_n - \mu_X) \xrightarrow{d} N(0, \sigma_X^2)$，其中 $\mu_X \neq 0$。

考虑函数 $g(y, x) = y/x$。在 $(\mu_Y, \mu_X)$ 处求偏导：
\[
\frac{\partial g}{\partial y} = \frac{1}{\mu_X}, \quad \frac{\partial g}{\partial x} = -\frac{\mu_Y}{\mu_X^2} .
\]

因此
\[
\sqrt{n}\left(\frac{Y_n}{X_n} - \frac{\mu_Y}{\mu_X}\right) \xrightarrow{d} N\left(0, \left(\frac{1}{\mu_X}, -\frac{\mu_Y}{\mu_X^2}\right) \begin{pmatrix} \sigma_Y^2 & \sigma_{YX} \\ \sigma_{YX} & \sigma_X^2 \end{pmatrix} \begin{pmatrix} \frac{1}{\mu_X} \\ -\frac{\mu_Y}{\mu_X^2} \end{pmatrix}\right) .
\]

展开即得渐近方差：
\[
\frac{\sigma_Y^2}{\mu_X^2} + \frac{\mu_Y^2\sigma_X^2}{\mu_X^4} - \frac{2\mu_Y\sigma_{YX}}{\mu_X^3} .
\tag*{\qed}
\]

% === 用户问答记录 ===
% (后续通过 interactive-qa 更新)
